% !TeX root = ../main.tex
\begin{frame}{01 研究背景——结构设计目标}
  \begin{SlideCanvas}
    \BodyAt{55}{84}{850}{支架的功能需求可转化为具体的\BICEmph{变形}、\BICEmph{承载}与\BICEmph{界面指标}，为超材料构型设计提供目标。}
    \PointAt{55}{145}{172}{功能需求}
    \PointAt{240}{145}{218}{工程问题}
    \PointAt{480}{145}{215}{结构设计量}
    \PointAt{720}{145}{180}{评价条件}
    \MainRuleAt{181}
    \PointAt{55}{200}{170}{病变覆盖}
    \DetailAt{240}{200}{220}{扩张时轴向缩短或伸长}
    \DetailAt{480}{200}{215}{周向与轴向变形耦合}
    \DetailAt{720}{200}{180}{长度及端部位移}
    \RuleAt{253}
    \PointAt{55}{272}{170}{管腔支撑}
    \DetailAt{240}{272}{220}{扩张阻力、回弹与径向抗压}
    \DetailAt{480}{272}{215}{非线性承载响应及材料行为}
    \DetailAt{720}{272}{180}{规定路径下的径向力与直径}
    \RuleAt{325}
    \PointAt{55}{344}{170}{贴壁与受力}
    \DetailAt{240}{344}{220}{接触间隙与局部应力集中}
    \DetailAt{480}{344}{215}{局部刚度与展开形状}
    \DetailAt{720}{344}{180}{血管接触模型}
    \RuleAt{397}
    \PointAt{55}{416}{170}{长期可靠性}
    \DetailAt{240}{416}{220}{周期变形与疲劳}
    \DetailAt{480}{416}{215}{杆件应变分布与柔顺性}
    \DetailAt{720}{416}{180}{循环载荷测试}
    \Sources{依据：\RefSCAI；器件评价例见\RefHas、\RefGupta。}
  \end{SlideCanvas}
  \BICConclusionBand{本次聚焦\BICBandEmph{长度控制}，并将径向承载作为同时满足的要求}
  \note{前两项构成本次研究主线。贴壁和组织应变进入接触评价，疲劳进入长期服役评价。结构优化可以缓和部分矛盾，但扩张不足还与钙化、尺寸选择和操作有关。需将结构收益与这些条件区分。径向力也不是组织接触状态的充分描述，见Borghi等，doi:10.1007/s11665-014-0913-z。}
\end{frame}
